% !TeX program = lualatex
% Motivationsschreiben Deutschlandstipendium – Erlind Sejdiu
% Kompilieren mit:  lualatex anschreiben.tex
% Layout bewusst identisch zum Lebenslauf (gleiche Schrift, gleicher Kopf).

\documentclass[a4paper,11pt]{article}

% ─── Pakete ──────────────────────────────────────────────────────────────────
\usepackage{fontspec}
\usepackage[ngerman]{babel}
\usepackage[left=2.4cm, right=2.4cm, top=1.8cm, bottom=1.6cm]{geometry}
\usepackage{microtype}
\usepackage{xcolor}
\usepackage[hidelinks]{hyperref}

\setmainfont{TeX Gyre Pagella}

\definecolor{ink}{HTML}{1C1C1C}
\definecolor{muted}{HTML}{6E6E6E}
\definecolor{hairline}{HTML}{C4C4C4}
\color{ink}

\hypersetup{pdfauthor={Erlind Sejdiu},
            pdftitle={Motivationsschreiben -- Deutschlandstipendium}}

\setlength{\parindent}{0pt}
\setlength{\parskip}{0.75em}
\pagestyle{empty}

% =============================================================================
\begin{document}

% ─── Kopf (identisch zum Lebenslauf) ─────────────────────────────────────────
{\LARGE Erlind Sejdiu}\\[0.35em]
{\small\color{muted}
Blücherstraße 18, Karlsruhe \enspace\textbar\enspace
0157\,3547\,4746 \enspace\textbar\enspace
\href{mailto:erlindamprogrammieren@gmail.com}{erlindamprogrammieren@gmail.com}}\\[0.45em]
\textcolor{hairline}{\rule{\linewidth}{0.7pt}}

\vspace{1.4em}

% ─── Empfänger ───────────────────────────────────────────────────────────────
{\parskip=0pt
Hochschule Karlsruhe\\
Stipendienkommission Deutschlandstipendium\\
Moltkestraße 30\\
76133 Karlsruhe\par}

\vspace{1.6em}

{\raggedleft\color{muted} Karlsruhe, \today \par}

\vspace{1.2em}

% TODO: Semester prüfen und ggf. anpassen
\textbf{Bewerbung um ein Deutschlandstipendium zum Wintersemester 2026/27}

\vspace{0.8em}

Sehr geehrte Damen und Herren,

seit März 2026 studiere ich Informatik im Master an der Hochschule Karlsruhe. Mein erstes Semester habe ich mit der Note 1,0 abgeschlossen. Bisher habe ich mein Studium komplett selbst finanziert. Ich arbeite als Werkstudent bei der ProSeS GmbH und betreibe nebenher eine GbR für Gartenarbeit, die ich mitgegründet habe. Das bringt mir viel Praxiserfahrung, kostet aber auch Zeit. Das Deutschlandstipendium würde mir helfen, etwas weniger zu arbeiten und mehr Zeit in mein Studium und eigene Projekte zu stecken.

Am meisten begeistert mich die Entwicklung nah an der Hardware, wo es wirklich auf Speicher und schnelle Laufzeiten ankommt. Meinen Bachelor in Technischer Informatik an der Hochschule Pforzheim habe ich mit der Note 1,3 gemacht. In meiner Bachelorarbeit mit der Note 1,0 bei ProSeS ging es um das Refactoring eines Benachrichtigungsmoduls. Im Praxissemester bei Vector Informatik habe ich an Software für einen Automotive Switch gearbeitet und bei Codewerk einen GUI Fuzzer für Cybersecurity entwickelt.

Neben dem Studium packe ich gerne mit an. Als Hilfskraft an der Hochschule habe ich anderen Studierenden bei Programmieraufgaben und Tests geholfen. In der Fachschaft und im Studierendenrat habe ich Events organisiert und mich für die Interessen der Studierenden eingesetzt. Auch die Arbeit in unserer GbR bringt mir viel. Dort habe ich gelernt, Verantwortung zu übernehmen, Kunden zu betreuen und mich auf mein Team zu verlassen.

Mit dem Stipendium könnte ich Arbeitszeit direkt in Zeit für Vorlesungen und eigene Projekte tauschen. Gerade bei Projekten an der Schnittstelle von Hardware und Software fehlen oft die Stunden. Über die Förderung an der Hochschule Karlsruhe würde ich mich sehr freuen.

\vspace{1.2em}

Mit freundlichen Grüßen

\vspace{2.0em}

Erlind Sejdiu

\end{document}
